\documentclass[11pt,a4paper]{article}

\usepackage[margin=2.5cm]{geometry}
\usepackage{amsmath,amssymb,amsthm,mathtools,bm}
\usepackage{booktabs,array}
\usepackage[hidelinks]{hyperref}
\usepackage[capitalise,nameinlink]{cleveref}

\newtheorem{theorem}{Theorem}[section]
\newtheorem{proposition}[theorem]{Proposition}
\newtheorem{lemma}[theorem]{Lemma}
\newtheorem{corollary}[theorem]{Corollary}
\theoremstyle{definition}
\newtheorem{definition}[theorem]{Definition}
\theoremstyle{remark}
\newtheorem{remark}[theorem]{Remark}

\newcommand{\R}{\mathbb{R}}
\newcommand{\dd}{\mathrm{d}}
\newcommand{\ind}{\mathbb{1}}
\newcommand{\what}{\widehat W}
\newcommand{\ee}{\bm e}
\newcommand{\nn}{\bm n}
\newcommand{\tv}{\bm t}
\newcommand{\xx}{\bm x}
\newcommand{\yy}{\bm y}
\newcommand{\sgn}{\operatorname{sgn}}
\newcommand{\atan}{\operatorname{atan2}}
\newcommand{\asinh}{\operatorname{asinh}}
\newcommand{\Lap}{\Delta}
\newcommand{\chord}{\mathrm{chord}}
\newcommand{\status}[1]{\hfill\textsf{[#1]}}

\title{Exact and asymptotic boundary integrals for SPH kernels:\\
divergence-theorem edge reductions, closed forms and curvature expansions}
\author{curvatureBoundaries project}
\date{\today}

\begin{document}
\maketitle

\begin{abstract}
We collect, in one logical chain, the mathematics behind the boundary integrals used to couple
smoothed-particle-hydrodynamics (SPH) fluids to solid boundaries: the kernel integral
$\lambda(\xx)=\int_S W(|\xx-\xx'|)\,\dd\xx'$ over a solid $S$, its gradient, its moments, and the
nodal weights of fields extended into the solid.  We prove (i) closed forms for planar and spherical
boundaries with Wendland and cubic-spline kernels, (ii) a divergence-theorem reduction of area
integrals over polygons to one-dimensional edge integrals, whose integrands are elementary for
piecewise-polynomial kernels, (iii) a compact-potential recursion for higher moments and its use for
FEM nodal weights, (iv) a closest-point curvature expansion and a slender-body (mean-value) series
for boundaries that are curved or thin compared to the support radius, and (v) the wall Laplacian
and wall-viscosity operators built on the same objects.  Every statement is tagged with how it was
checked (symbolically in Maple, against independent quadrature, or not at all), and an appendix lists
the problems found while auditing the existing note collection.
\end{abstract}

\tableofcontents

\section{Setting and conventions}\label{sec:setting}

\subsection{Kernels}
Throughout, $n\in\{2,3\}$ is the spatial dimension and $h$ the \emph{support radius}. A kernel is
\begin{equation}\label{eq:kernel}
  W(r;h)=\frac{C_n}{h^n}\,\what\!\left(\frac rh\right),\qquad \what\in C([0,1]),\quad \what(q)=0\ (q\ge1),
\end{equation}
with $C_n$ fixed by $\int_{\R^n}W=1$, i.e.\ $2\pi C_2\int_0^1\what\,q\,\dd q=1$ and
$4\pi C_3\int_0^1\what\,q^2\dd q=1$. We work at $h=1$ and restore $h$ by scaling (\cref{lem:scaling}).
The kernels used are listed in \cref{tab:kernels}; all are piecewise polynomial in $q$.
\begin{table}[h]
\centering
\begin{tabular}{llll}
\toprule
kernel & $\what(q)$ & $C_2$ & $C_3$\\
\midrule
cubic & $(1-q)^3-4(\tfrac12-q)_+^3$ & $\frac{80}{7\pi}$ & $\frac{16}{\pi}$\\[2pt]
Wendland C2 (w2) & $(1-q)^4(1+4q)$ & $\frac7\pi$ & $\frac{21}{2\pi}$\\[2pt]
Wendland C4 (w4) & $(1-q)^6(1+6q+\frac{35}{3}q^2)$ & $\frac9\pi$ & $\frac{495}{32\pi}$\\[2pt]
Wendland C6 (w6) & $(1-q)^8(1+8q+25q^2+32q^3)$ & $\frac{78}{7\pi}$ & $\frac{1365}{64\pi}$\\
\bottomrule
\end{tabular}
\caption{Kernels at $h=1$ (Wendland~\cite{wendland1995}). On $[0,\frac12]$ the cubic equals
$\frac12-3q^2+3q^3$, so $\what(0)=\frac12$.}\label{tab:kernels}
\end{table}

\begin{lemma}[Normalisation]\label{lem:norm}
The constants in \cref{tab:kernels} satisfy $\int_{\R^n}W=1$ for $n=2,3$.
\end{lemma}
\begin{proof}
Direct integration of polynomials; checked symbolically (Maple, \texttt{20\_paper\_audit.mpl}).
\end{proof}

\subsection{The boundary integrals}
Let $S\subset\R^n$ be the solid and $\xx$ the evaluation point (a fluid particle). With
$\yy=\xx'-\xx$ and $r=|\yy|$ we study
\begin{align}
  \lambda(\xx)&=\int_S W(|\xx-\xx'|)\,\dd\xx', \label{eq:lambda}\\
  m_{\alpha}(\xx)&=\int_S \yy^\alpha\,W(r)\,\dd\xx',\qquad |\alpha|=k, \label{eq:moments}\\
  g_{\alpha}(\xx)&=\int_S \yy^\alpha\,\nabla_{\xx}W(|\xx-\xx'|)\,\dd\xx'. \label{eq:galpha}
\end{align}
The value $\lambda$ is the volume of solid seen by the particle; $m_\alpha$ and $g_\alpha$ are needed
for fields extended into the solid (\cref{sec:fem}). The distance from $\xx$ to $\partial S$ is $d\ge0$
for a fluid-side particle ($d<0$ inside the solid, $\lambda(-d)=1-\lambda(d)$ for a flat wall).

\begin{lemma}[Scaling]\label{lem:scaling}
$\lambda(d;h)=\lambda(d/h;1)$ and $m_\alpha(d;h)=h^{|\alpha|}\,m_\alpha(d/h;1)$, where $d$ stands for the whole
geometry measured in the same unit as $h$. Consequently $\nabla_{\xx}\lambda$ carries a factor $h^{-1}$ and
$\partial_{\xx}m_\alpha$ a factor $h^{|\alpha|-1}$.
\end{lemma}
\begin{proof}
Substitute $\xx'=\xx+h\yy_1$ in \cref{eq:lambda,eq:moments}: $\dd\xx'=h^n\dd\yy_1$ cancels the $h^{-n}$ in
\cref{eq:kernel}, and $\yy^\alpha=h^{|\alpha|}\yy_1^\alpha$.
\end{proof}

\subsection{The idea in one page}\label{sec:idea}
Everything below is one trick used three times: \emph{write the integrand as a divergence, so the area (volume)
integral over the solid becomes a boundary integral, and only the part of the boundary inside the support radius
contributes.}
\begin{center}\small
\begin{tabular}{>{\raggedright\arraybackslash}p{3.3cm}>{\raggedright\arraybackslash}p{6.2cm}>{\raggedright\arraybackslash}p{3cm}}
\toprule
quantity & vector field whose divergence (or gradient) is the integrand & result\\
\midrule
value $\int_T g$ & $\bm G=\yy\,M_g(r)/r^{2}$ with $\nabla\!\cdot\bm G=g$ & \cref{thm:value}\\
gradient $\nabla_{\xx}\!\int_Tg$ & translate the domain: $\partial_\delta\!\int_{T-\delta\ee}g=-\oint n\,g$ & \cref{thm:grad}\\
moments $\int_T\yy^\alpha f$ & $\yy^\beta\,\Phi[f]$ with $\nabla\Phi=\yy f$ & \cref{thm:recursion}\\
\bottomrule
\end{tabular}
\end{center}
For a polygon the boundary is a union of edges; on each edge the integrand is an explicit function of the single
variable $s$ (position along the edge), and for piecewise-polynomial kernels these one-dimensional integrals are
elementary (\cref{prop:primitives}).  Sections~\ref{sec:planar}--\ref{sec:edge} give the exact oracle and the edge calculus;
\cref{sec:fem,sec:closed} use them for fields on elements and for closed boundaries; \cref{sec:curv,sec:slender}
treat smooth or thin bodies asymptotically instead of adding edges; \cref{sec:threed,sec:wall} extend and use the
results.

\subsection{Notation}\label{sec:notation}
\begin{center}\small
\begin{tabular}{lp{11.4cm}}
\toprule
symbol & meaning\\
\midrule
$\xx$, $\xx'$, $\yy$, $r$ & evaluation point, integration point, $\yy=\xx'-\xx$, $r=|\yy|$\\
$h$, $q=r/h$ & support radius of the kernel, scaled distance ($h=1$ in all derivations)\\
$d$, $z_e$ & distance of $\xx$ to a wall (\cref{sec:planar}); signed distance of $\xx$ to the line of edge $e$ (\cref{sec:edge}), $z_e=\nn_e\cdot(p_e-\xx)$\\
$r_c$ & \emph{cut radius} of a radial profile $g$ ($g=0$ for $r>r_c$); $r_c=1$ for a full kernel, $r_c=\frac12$ for the inner cubic block\\
$R$, $\kappa=1/R$ & radius and curvature of a disk or ball (\cref{sec:sphere,sec:closed,sec:curv}); never a cut radius\\
$M_g(r)$ & $\int_0^rt\,g(t)\,\dd t$ (2D) or $\int_0^rt^2g(t)\,\dd t$ (3D), the \emph{radial mass} of $g$\\
$\Phi[f]$, $\varphi(\sigma)$ & compact potential of a radial profile $f$ (\cref{thm:recursion}); the kernel as a function of $\sigma=r^2$ (\cref{sec:curv})\\
$I_m(s)$ & $\int_0^sr^m\dd s$ along an edge line (\cref{prop:primitives}); not to be confused with the
planar family $\mathcal J_m(d)$ of \cref{lem:J}\\
$\lambda_\Delta$ & $\int_S\nabla^2_{\xx}W$, the wall Laplacian (\cref{sec:wall}); a number, not a difference\\
$\ind[\cdot]$ & indicator\\
\bottomrule
\end{tabular}
\end{center}

\subsection{Verification tags}
Each result carries a tag: \textsf{[M]}~checked symbolically in Maple (exact arithmetic, or 40--50 digit
quadrature where stated); \textsf{[N]}~checked numerically against an independent quadrature
(\texttt{tests/}); \textsf{[A]}~argued here but not independently checked; \textsf{[P]}~plan, not proved.
The scripts are \texttt{maple/10--14} (project) and \texttt{maple/20\_paper\_audit.mpl} (this paper).

\section{Planar and spherical boundaries}\label{sec:planar}

Set $h=1$. Place the particle at the origin and let the boundary be the plane (line) at distance
$d\in[0,1]$, the solid being the far side. $\lambda$ then depends on $d$ alone.

\subsection{Three dimensions}
\begin{proposition}[Shell form]\label{prop:planar3}
$\displaystyle \lambda_3(d)=\int_d^1 C_3\what(q)\,2\pi q(q-d)\,\dd q.$ \status{M}
\end{proposition}
\begin{proof}
The sphere of radius $q$ about the particle meets the half space $z\ge d$ in a cap of height $q-d$
(for $q\ge d$) and area $2\pi q(q-d)$ (Archimedes); integrate shell by shell.
\end{proof}

\paragraph{Wendland kernels.} With $\what=(1-q)^m P(q)$, $P=\sum_p c_pq^p$, put
$F(m,p,x)=\int_0^x t^p(1-t)^m\dd t=\sum_{j=0}^m(-1)^j\binom mj\frac{x^{p+j+1}}{p+j+1}$. Then
\begin{equation}\label{eq:lam3-wend}
 \lambda_3(d)=2\pi C_3\sum_p c_p\Big[F(m,p{+}2,1)-F(m,p{+}2,d)-d\big(F(m,p{+}1,1)-F(m,p{+}1,d)\big)\Big],
\end{equation}
a single polynomial in $d$ of degree $m+\deg P+3$, i.e.\ $8,11,14$ for w2, w4, w6, with $\lambda_3(0)=\frac12$,
$\lambda_3(1)=0$. \status{M}

\paragraph{Cubic spline.} The knot at $q=\frac12$ splits the integral (inner correction $-4(\frac12-q)^3$
only for $q\le\frac12$):
\begin{align}
 d\in[0,\tfrac12]:&\quad \lambda_3=\tfrac1{60}\,(192d^6-288d^5+160d^3-84d+30),\\
 d\in[\tfrac12,1]:&\quad \lambda_3=-\tfrac{8}{15}\,(2d^6-9d^5+15d^4-10d^3+3d-1),
\end{align}
which agree at $d=\frac12$ (value $\frac1{30}$). \status{M} (Winchenbach et al.~\cite{winchenbach2020} give the cubic case.)

\subsection{Two dimensions}\label{sec:planar2}
\begin{proposition}[Arc form]\label{prop:planar2}
$\displaystyle\lambda_2(d)=\int_d^1 C_2\what(q)\,2q\arccos\!\big(\tfrac dq\big)\,\dd q.$ \status{M}
\end{proposition}
\begin{proof}
The circle of radius $q\ge d$ about the particle meets the half plane $y\ge d$ where $q\cos\phi\ge d$, i.e.\ in the arc
$|\phi|\le\arccos(d/q)$ of length $2q\arccos(d/q)$. Integrate shell by shell. (Compare 3D: there the shell area is
$2\pi q\,(q-d)$, polynomial in $q$; here the inverse cosine makes the integrand non-polynomial.)
\end{proof}

\begin{proposition}[By-parts reduction]\label{prop:byparts}
Let $v_0(q)=\int_0^q t\,\what(t)\,\dd t$, so that $2\pi C_2v_0(1)=1$. Then
\begin{equation}\label{eq:lam2-byparts}
 \lambda_2(d)=2C_2\Big[v_0(1)\arccos d-d\int_d^1\frac{v_0(q)}{q\sqrt{q^2-d^2}}\,\dd q\Big].
\end{equation}
In particular $\lambda_2(0)=\frac12$, $\lambda_2(1)=0$ and
$\lambda_2'(d)=-2C_2\int_d^1 \what(q)\,q\,(q^2-d^2)^{-1/2}\dd q$. \status{M}
\end{proposition}
\begin{proof}
Integrate \cref{prop:planar2} by parts with $u=\arccos(d/q)$ and $\dd v=2C_2q\,\what(q)\,\dd q$, i.e.\ $v=2C_2v_0$.
Since $\frac{\dd}{\dd q}\arccos(d/q)=d/(q\sqrt{q^2-d^2})$, $u(d)=\arccos1=0$ and $u(1)=\arccos d$,
\[
 \lambda_2=\big[u\,v\big]_d^1-\int_d^1v\,u'\,\dd q
 =2C_2v_0(1)\arccos d-2C_2\,d\int_d^1\frac{v_0(q)}{q\sqrt{q^2-d^2}}\dd q .
\]
At $d=0$ the second term vanishes and $\lambda_2(0)=2C_2v_0(1)\cdot\frac\pi2=\frac12$. For the derivative
differentiate the arc form directly: the boundary term at $q=d$ vanishes because $\arccos1=0$, and
$\partial_d\arccos(d/q)=-(q^2-d^2)^{-1/2}$.
\end{proof}

\begin{lemma}[The family $\mathcal J_m$]\label{lem:J}
Let $\mathcal J_m(d)=\int_d^1 q^m(q^2-d^2)^{-1/2}\dd q$. Then
$\mathcal J_0=\ln\frac{1+\sqrt{1-d^2}}{d}$, $\mathcal J_1=\sqrt{1-d^2}$ and
\begin{equation}\label{eq:Jrec}
 m\,\mathcal J_m=\sqrt{1-d^2}+(m-1)\,d^2\mathcal J_{m-2}\qquad(m\ge2).
\end{equation}
Consequently, for odd $m$, $\mathcal J_m=\sqrt{1-d^2}\,\mathrm{poly}(d^2)$; for even $m$,
$\mathcal J_m=\sqrt{1-d^2}\,\mathrm{poly}(d^2)+\mathrm{poly}(d^2)\,\mathcal J_0$, where $\mathcal J_0$ contains
$\ln d$ and $\ln(1+\sqrt{1-d^2})$. \status{M}
\end{lemma}
\begin{proof}
$\mathcal J_0=\operatorname{arcosh}(1/d)$ and $\mathcal J_1$ are direct. For $m\ge2$, using
$\sqrt{q^2-d^2}=(q^2-d^2)/\sqrt{q^2-d^2}$,
\[
 \frac{\dd}{\dd q}\Big(q^{m-1}\sqrt{q^2-d^2}\Big)
 =(m-1)q^{m-2}\sqrt{q^2-d^2}+\frac{q^m}{\sqrt{q^2-d^2}}
 =\frac{m\,q^m-(m-1)d^2q^{m-2}}{\sqrt{q^2-d^2}} .
\]
Integrate over $[d,1]$: the left side gives $\sqrt{1-d^2}$ (the term at $q=d$ vanishes). The parity statements follow
by induction from $\mathcal J_1$ and $\mathcal J_0$. For example $\mathcal J_3=\frac13(1+2d^2)\sqrt{1-d^2}$ and
$\mathcal J_2=\frac12(\sqrt{1-d^2}+d^2\mathcal J_0)$.
\end{proof}
\begin{remark}
The project notes state that odd $\mathcal J_m$ lie in $\mathrm{poly}(d^2)\{1,\sqrt{1-d^2}\}$; \cref{eq:Jrec}
shows the constant is never present (\cref{app:audit}).
\end{remark}

\paragraph{Wendland.} Because $\what(q)=1+O(q^2)$, $v_0=O(q^2)$ is a polynomial and $v_0(q)/q=\sum_{m\ge1}a_mq^m$. Then
\eqref{eq:lam2-byparts} reads
\begin{equation}\label{eq:lam2-wend}
 \lambda_2(d)=2C_2\Big[v_0(1)\arccos d-d\sum_{m\ge1}a_m\mathcal J_m(d)\Big],
\end{equation}
that is $\arccos d$ plus $d\,[\mathrm{poly}\sqrt{1-d^2}+\mathrm{poly}\ln d+\mathrm{poly}\ln(1+\sqrt{1-d^2})]$.
\emph{Example (w2).} $\what=1-10q^2+20q^3-15q^4+4q^5$ gives
$v_0=\frac{q^2}2-\frac{5q^4}2+4q^5-\frac{5q^6}2+\frac{4q^7}7$, $v_0(1)=\frac1{14}$ (so $2C_2v_0(1)=\frac1\pi$ with $C_2=\frac7\pi$), and
\[
 \lambda_2^{\rm w2}(d)=\frac{14}\pi\Big[\frac{\arccos d}{14}-d\Big(\tfrac12\mathcal J_1-\tfrac52\mathcal J_3+4\mathcal J_4-\tfrac52\mathcal J_5+\tfrac47\mathcal J_6\Big)\Big]
 \qquad\text{[M, 40 digits at }d=\tfrac3{10}].
\]
\paragraph{Cubic.} The knot splits the range. For $d\ge\frac12$ only the outer piece $\what=(1-q)^3$ is present and the
Wendland recipe applies. For $d<\frac12$ the inner correction $-4(\frac12-q)^3$ is a second block with upper limit
$\frac12$; its integrals end at $q=\frac12$ instead of $1$, which adds terms in $\arccos2d$, $\sqrt{1-4d^2}$ and
$\ln(1+\sqrt{1-4d^2})$. At $d=\frac12$ only the outer piece contributes, and
\[
 \lambda_2(\tfrac12)=\frac{64\pi-294\sqrt3+249\ln(2+\sqrt3)}{168\pi}\approx0.03744 \quad\text{[M, 40 digits]}.
\]

\subsection{Solid ball in 3D}\label{sec:sphere}
Solid ball of radius $R$, particle outside at distance $d$ from the surface, hence $D=R+d$ from the centre.
\begin{proposition}[Cap area]\label{prop:sphere}
The part of the sphere $|\yy|=q$ (centre at the particle) that lies inside the ball has area
$A(q)=\dfrac{\pi q}{D}\,(R^2-D^2+2Dq-q^2)$ for $q\in(D-R,D+R)$, and $0$ otherwise.
Hence $\lambda_{\rm sph}(R,d)=\int_d^{\min(1,\,2R+d)}C_3\what(q)A(q)\,\dd q$. \status{M}
\end{proposition}
\begin{proof}
A point of the sphere at polar angle $\theta$ from the direction to the ball centre has squared distance
$q^2+D^2-2Dq\cos\theta$ to that centre, so it lies in the ball iff $\cos\theta\ge(q^2+D^2-R^2)/(2Dq)$. This is a spherical
cap of height $h_c=q(1-\cos\theta_0)=q-(q^2+D^2-R^2)/(2D)$ and area $2\pi q\,h_c$ (Archimedes), which is the stated
$A$. It vanishes at $q=D\mp R$ [M], where the sphere touches the ball from outside or from inside. Beyond $q=D+R$ the
sphere encloses the ball, so $A=0$.
\end{proof}
Two branches arise: $2R+d\ge1$ (the support truncates, upper limit $1$) and $2R+d<1$ (the ball lies entirely inside the support,
upper limit $2R+d$). For Wendland kernels $A$ is quadratic in $q$, so the $F$-basis of \cref{eq:lam3-wend} closes both
branches (with one more power of $q$) and $\lambda_{\rm sph}=N(R,d)/(R+d)$ with $N$ a polynomial; on branch~1, $N$ is linear in $R$ because
$R^2-D^2=-(2Rd+d^2)$ is. Since $\lim_{R\to\infty}A=2\pi q(q-d)$ \status{M}, branch~1 recovers
\cref{eq:lam3-wend}.

\subsection{Consistency of the two routes}
Planar closed forms are the exact oracles against which every later construction is compared. Their
independence from the edge calculus of \cref{sec:edge} is established in \cref{prop:halfplane}.

\section{Edge calculus in two dimensions}\label{sec:edge}

This is the core of the construction. For a polygon $T$ we reduce $\int_T W$, its gradient and all
moments to one-dimensional integrals along the edges. For piecewise-polynomial $W$ these are elementary.
The plan is the table of \cref{sec:idea}: a value identity (\cref{thm:value}), a gradient identity (\cref{thm:grad})
and a recursion for moments (\cref{thm:recursion}); \cref{prop:primitives} supplies the one-dimensional integrals they need.

\subsection{Geometry}
Let $T$ be a simple polygon, counter-clockwise. For the evaluation point $\xx$ and an edge $e=(p_e,q_e)$ let
$\tv_e$ be the unit tangent, $\nn_e$ the \emph{outward} normal, and
\begin{gather}
 z_e=\nn_e\!\cdot(p_e-\xx),\qquad s_0=\tv_e\!\cdot(p_e-\xx),\qquad s_1=\tv_e\!\cdot(q_e-\xx),\notag\\
 \yy=z_e\nn_e+s\,\tv_e,\qquad r=\sqrt{s^2+z_e^2}\quad\text{on the edge line.}
\end{gather}
So $s$ is the position along the edge measured from the foot of the perpendicular from $\xx$, and $z_e$ is the distance to the edge line:
$z_e>0$ iff $\xx$ lies on the inner side of the edge line. For a profile with cut radius $r_c$, the
edge meets the support in the \emph{chord} $s\in[\max(s_0,-\ell),\min(s_1,\ell)]$, $\ell=\sqrt{r_c^2-z_e^2}$, empty if
$|z_e|\ge r_c$ or the interval is empty.

\subsection{Elementary primitives}
With $r=\sqrt{s^2+z^2}$ define $I_m(s)=\int_0^s r^m\,\dd s$ and $S_{j,m}=\int s^jr^m\dd s$.

\begin{proposition}[Primitives]\label{prop:primitives}
(a) $(m+1)I_m=s\,r^m+m z^2I_{m-2}$ with $I_0=s$, $I_{-1}=\asinh(s/|z|)$, $zI_{-2}=\arctan(s/z)$.
Hence for even $m=2k\ge0$, $I_m=\sum_i\binom ki z^{2k-2i}\frac{s^{2i+1}}{2i+1}$ (a polynomial), and for odd $m\ge1$,
$I_m=\frac{m!!}{(m+1)!!}z^{m+1}\asinh(s/|z|)+s\,r\,\mathrm{poly}(r^2,z^2)$.\\
(b) $\int s\,r^m\dd s=r^{m+2}/(m+2)$ for $m\neq-2$, and $\ln r$ for $m=-2$.\\
(c) Powers of $s$ are reduced with $s^2=r^2-z^2$. For $j=2p$ even,
$S_{j,m}=\sum_{i=0}^{p}\binom pi(-z^2)^{p-i}I_{m+2i}$; for $j=2p+1$ odd,
$S_{j,m}=\sum_{i=0}^{p}\binom pi(-z^2)^{p-i}\,\dfrac{r^{m+2i+2}}{m+2i+2}$ (with $\ln r$ if $m+2i+2=0$). Note that $p=\lfloor j/2\rfloor$ in both cases.\\
(d) For a chord $[lo,hi]$ and either sign of $z$:
$\arctan(hi/z)-\arctan(lo/z)=\atan\!\big(z(hi-lo),\,z^2+hi\cdot lo\big)$. \status{M, maple/10}
\end{proposition}
\begin{proof}
(a) Differentiating $s\,r^m$ gives $r^m+m s^2r^{m-2}=(m+1)r^m-mz^2r^{m-2}$, which is the recurrence; the closed forms follow by
induction from $I_0$ and $I_{1}=\frac12\big(s\,r+z^2\asinh(s/|z|)\big)$. (b) is the chain rule. (c) For $j=2p$ expand
$s^{2p}r^m=(r^2-z^2)^pr^m$ binomially and integrate term by term; for $j=2p+1$ write $s^{2p+1}r^m=s\,(r^2-z^2)^p r^m$ and use (b).
(d) is the tangent subtraction formula $\tan(a-b)=\frac{\tan a-\tan b}{1+\tan a\tan b}$ with the denominator and numerator multiplied
by $z^2>0$; the $\atan$ form selects the correct branch since the difference lies in $(-\pi,\pi)$ and has the sign of $z(hi-lo)$.
All derivative identities were verified symbolically for $|m|\le12$, $j\le7$.
\end{proof}
\begin{remark}
$I_m$ is odd in $s$, but $\int s\,r^m\dd s$ is \emph{even}; a chord spanning $s=0$ therefore adds $I_m$ contributions but
subtracts equal contributions of the second kind. For $m<-2$ the downward recurrence divides by $z^2$ each step and is numerically
unusable at small $z$; only $m\ge0$ is needed by the production path.
\end{remark}

\subsection{Truncated monomials}
Any piecewise-polynomial radial kernel is a finite combination of \emph{blocks} $f_{n,r_c}(r)=r^n\ind[r\le r_c]$:
$W=\sum_jq_j(r)\ind[r\le r_j]$ with $q_{\rm last}$ the outer piece and $q_j=\text{piece}_j-\text{piece}_{j+1}$.
\begin{itemize}
\item Wendland: one block family with $r_c=1$.
\item Cubic: $W=C_2[(1-q)^3-4(\tfrac12-q)^3\ind[q\le\frac12]]$, blocks at $r_c=1$ and $r_c=\tfrac12$.
\end{itemize}
\status{M}

\subsection{Value identity}
For a radial $g$ supported on $r\le r_c$ set $M_g(r)=\int_0^r t\,g(t)\,\dd t$ (so $M_g(r)=M_g(r_c)$ for $r\ge r_c$).
\begin{theorem}[Value identity]\label{thm:value}
Let $\xx\notin\partial T$. Then
\begin{equation}\label{eq:value}
 \int_T g(|\xx'-\xx|)\,\dd\xx'=\ind[\xx\in T]\,2\pi M_g(r_c)+\sum_{e\in\partial T}z_e\int_{\chord_e}\frac{M_g(r)-M_g(r_c)}{r^2}\,\dd s.
\end{equation}
For a normalised kernel $2\pi M_W(1)=1$, so the first term is $\ind[\xx\in T]$.
\end{theorem}
\begin{proof}
\emph{Step 1: $g$ is a divergence.} Put $\bm G(\yy)=\yy\,M_g(r)/r^2$. Away from $\yy=0$,
$\nabla\!\cdot\bm G=\dfrac{2M_g}{r^2}+r\,\dfrac{\dd}{\dd r}\dfrac{M_g}{r^2}=\dfrac{2M_g}{r^2}+\dfrac{M_g'}{r}-\dfrac{2M_g}{r^2}=\dfrac{M_g'}r=g$, since $M_g'=r\,g$.\\
\emph{Step 2: divergence theorem.} If $\xx\notin T$, $\bm G$ is smooth on $T$ and $\int_Tg=\oint_{\partial T}\bm G\cdot\nn\,\dd s$.
If $\xx\in T$, apply it on $T\setminus B_\varepsilon(\xx)$; the flux through the small circle (normal pointing towards $\xx$) is
$-\frac{M_g(\varepsilon)}{\varepsilon}\cdot2\pi\varepsilon=-2\pi M_g(\varepsilon)\to0$ since $M_g(\varepsilon)=O(\varepsilon^2)$.
On an edge $\nn_e\!\cdot\!\yy=z_e$, so in both cases
\begin{equation}\label{eq:value-full}
 \int_Tg=\sum_e z_e\int_e\frac{M_g(r)}{r^2}\,\dd s\quad(\text{whole edges}).
\end{equation}
\emph{Step 3: localise.} Write $M_g=\big(M_g-M_g(r_c)\big)+M_g(r_c)$. The first part vanishes for $r\ge r_c$, so only the chord contributes.
For the constant part, $z_e\!\int\!\dd s/r^2=\arctan(s/z_e)\big|_{s_0}^{s_1}$ is the angle under which edge $e$ is seen from $\xx$, and
$\sum_e$ of these is the total angle subtended by the closed curve $\partial T$, namely $2\pi$ if $\xx\in T$ and $0$ otherwise.
\end{proof}
This is the identity behind all edge-based boundary methods; the 3D analogue replaces $2\pi$ by $4\pi$ and
the angle by the solid angle (\cref{sec:threed}). In words: only the edges near $\xx$ matter, plus a topological term that
counts whether $\xx$ is inside.

\paragraph{Block form.} For $g=r^n\ind[r\le r_c]$ one has $M_g=\min(r,r_c)^{n+2}/(n+2)$ and
$(M_g-M_g(r_c))/r^2=(r^n-r_c^{n+2}r^{-2})/(n+2)$ on the chord, so by $z\!\int\!\dd s/r^2=\arctan(s/z)$,
\begin{equation}\label{eq:block}
 \int_T r^n\ind[r\le r_c]=\ind[\xx\in T]\frac{2\pi r_c^{n+2}}{n+2}+\sum_e\frac{1}{n+2}\Big[z_e\big(I_n\big|_{lo}^{hi}\big)-r_c^{n+2}\,\Delta\!\arctan\Big].
\end{equation}
For a multi-block kernel apply \cref{eq:block} to each block and sum. Collecting the blocks of cut radius $r_j$ into one
group with radial mass $M_j(r_j)$, the group contributes the indicator weight $w_j=2\pi M_j(r_j)$ and the term
$-\frac{w_j}{2\pi}\Delta\!\arctan_j$ (the angle taken on the chord of radius $r_j$). For the cubic spline $w_j=\frac87$ ($r_j=1$) and $-\frac17$
($r_j=\frac12$); the \emph{negative} weight is correct and $\sum_jw_j=1$. \status{M}

\subsection{Degenerate placements}\label{sec:degenerate}
$\lambda_T(\xx)=\int_TW$ is continuous in $\xx$ ($W$ is bounded), but the two parts of \eqref{eq:value} are not.
Take $\xx$ crossing the line of an edge $e$ at a point projecting into the edge interior, from outside to inside, so $z_e$ changes sign from $-$ to $+$.
The indicator jumps by $+2\pi M_W(1)=+1$. In the chord term of $e$ the integrand $z_e\,(M_W(r)-M_W(1))/r^2$ concentrates at $s=0$ as $|z_e|\to0$,
where $M_W(0)=0$: it tends to $-M_W(1)\,\pi\,\sgn z_e$, i.e.\ it jumps by $-2\pi M_W(1)=-1$. The two jumps cancel; all other edges are continuous.
Therefore the value \emph{on} the boundary, obtained by continuity, uses the average of the two one-sided limits of each part:
\emph{$\xx$ in the open interior of an edge: indicator $\frac12$, that edge contributes $0$ to the chord terms (and to every arctan of \cref{eq:block});
$\xx$ at a vertex: indicator = interior angle$/2\pi$; an edge with $z_e=0$ contributes $0$.}
(At a vertex, the two incident edges each contribute $-M_W(1)\,(\pi-\frac\alpha2)$ in the limit from the bisector, for a total $-M_W(1)(2\pi-\alpha)$,
and $2\pi M_W(1)-M_W(1)(2\pi-\alpha)=M_W(1)\,\alpha$ is exactly the indicator $\alpha/2\pi$ times $2\pi M_W(1)$.)
These conventions must use the \emph{same} orientation predicate for the sign of $z_e$ and for the indicator. \status{N}

\paragraph{Worked check: a flat wall.} For a half-plane solid at distance $d$ (single edge, $z=-d$, indicator $0$), \cref{eq:value} gives
$\lambda(d)=-d\int_{-\ell}^{\ell}\frac{M_W(r)-M_W(1)}{r^2}\dd s$. As $d\to0$ the factor $d$ kills the $M_W(r)\sim r^2$ part
(the integrand is bounded) and what remains is $+d\,M_W(1)\int_{-\ell}^{\ell}\dd s/r^2=2M_W(1)\arctan(\ell/d)\to\pi M_W(1)=\frac12$,
the known value $\lambda(0)=\frac12$ of \cref{prop:byparts}.

\subsection{Gradient identity}
\begin{theorem}[Gradient identity]\label{thm:grad}
For $g$ bounded, piecewise continuous with jumps on circles about $\xx$ that meet $\partial T$ transversally,
\[
 \partial_{x_j}\int_Tg(|\xx'-\xx|)\,\dd\xx'=-\sum_e n_{e,j}\int_{\chord_e}g(r)\,\dd s.
\]
In particular, $\nabla_{\xx}\lambda=-\sum_e\nn_e\int_{\chord_e}W\,\dd s$ \emph{without any indicator term}, and the identity holds
block by block.
\end{theorem}
\begin{proof}
$\int_Tg(\xx'-\xx)\dd\xx'=\int_{T-\xx}g(\yy)\dd\yy$, so moving $\xx$ by $\delta\ee_j$ translates the domain by $-\delta\ee_j$.
A domain translated by $\bm v$ gains $\oint(\bm v\!\cdot\!\nn)\,g\,\dd s+o(|\bm v|)$ (Reynolds transport; $g$ is continuous across
$\partial T-\xx$ a.e.), here $-\delta\oint n_jg\,\dd s$. Only the chords carry $g\neq0$.
\end{proof}
Sign check: moving $\xx$ towards the solid, i.e.\ against $\nn_e$ (which points out of the solid), increases $\lambda$; the formula gives $\nabla\lambda\parallel-\nn_e$.
\begin{corollary}[First moment]\label{cor:firstmoment}
With $\Psi(r)=\int_r^1tW\dd t$ (so $\yy W=-\nabla\Psi$),
\[
 \int_T\yy W\dd\xx'=-\sum_e\nn_e\int_{\chord_e}\Psi\,\dd s. \qquad\status{N}
\]
This is the case $|\alpha|=1$ of \cref{thm:recursion} with $\Psi=-\Phi[W]$.
\end{corollary}
\begin{corollary}\label{cor:fullspace}
Gradients of moments are edge-local:
$\partial_{x_j}m_\alpha=-\sum_en_{e,j}\int_{\chord_e}\yy^\alpha W\dd s$. (The proof of \cref{thm:grad} translates the domain $T-\xx$
and so differentiates \emph{both} $\yy^\alpha$ and $W$.) If $T$ covers the support the sum
vanishes, as it must.
\end{corollary}

\subsection{Compact-potential recursion for moments}
Let $f$ be radial and supported on $r\le r_c$ and $\Phi[f](r)=-\int_r^{r_c}t\,f(t)\,\dd t=M_f(r)-M_f(r_c)$. Then $\Phi$ is compactly
supported, continuous, and $\partial_i\Phi=y_if$ (since $\Phi'(r)=rf$ and $\partial_i r=y_i/r$).
\begin{theorem}[Recursion]\label{thm:recursion}
For $\alpha=\beta+\ee_i$,
\begin{equation}\label{eq:recursion}
 \int_T\yy^{\alpha}f\,\dd\xx'=\sum_en_{e,i}\int_{\chord_e}\yy^{\beta}\Phi\,\dd s-\beta_i\int_T\yy^{\beta-\ee_i}\Phi\,\dd\xx'.
\end{equation}
\end{theorem}
\begin{proof}
$\partial_i(\yy^\beta\Phi)=\beta_i\yy^{\beta-\ee_i}\Phi+\yy^{\beta}y_if$; integrate over $T$ and use the divergence theorem on the left.
$\Phi$ is Lipschitz, so no jump term arises even if $f$ jumps at $r_c$. \status{M, $k\le4$}
\end{proof}
The degree drops by $2$ per step. For odd $k$ the recursion ends in pure edge integrals. For even $k$ it ends in a
value-type integral of an iterated potential, evaluated by \cref{thm:value} with $M_g$ of that potential --- which is why \cref{thm:value}
is stated for general $g$ with $2\pi M_g(r_c)\neq1$. For a block $f=r^n\ind[r\le r_c]$, $\Phi=\frac{r^{n+2}-r_c^{n+2}}{n+2}\ind[r\le r_c]$ again lies in the
block family, so only $I_m$ and $\int s\,r^m\dd s$ with $m\ge0$ occur.

\emph{Numerical confirmation.} For a non-convex hexagon, all four kernels, particle inside, outside and near an edge, all $28$ multi-indices with $|\alpha|\le6$
agree with an independent ray quadrature of the area integral to relative $3\cdot10^{-11}$ (\path{scripts/derivation_checks/paper_edge_identities_k6.py}); this covers the range $k\le6$ needed by \cref{sec:curv}
beyond the symbolic check. The edge form of $\partial_{\xx}m_\alpha$ (\cref{cor:fullspace}) agrees with central differences to the finite-difference error ($\le10^{-5}$) and
$g_\alpha=\partial_{\xx}m_\alpha+\sum_j\alpha_j\ee_jm_{\alpha-\ee_j}$ (\cref{prop:fem}) with a direct ray quadrature of $\yy^\alpha\nabla_{\xx}W$ to $4\cdot10^{-12}$. \status{N}

\emph{Example ($k=2$).} For $\alpha=2\ee_1$ take $\beta=\ee_1$, $i=1$:
$\int_Ty_1^2f=\sum_en_{e,1}\int_{\chord_e}y_1\Phi\,\dd s-\int_T\Phi$, where on the chord $y_1=z_en_{e,1}+s\,t_{e,1}$ is linear in $s$, and
$\int_T\Phi$ is a value integral with $M_\Phi(r_c)\ne\frac1{2\pi}$ in general.

\begin{remark}[Direct form (b) fails on an edge line]
An alternative with $\ind[\xx\in T]\oint\omega^\alpha r_c^{n+2+k}/(n+2+k)+\sum_e\frac{z_e}{n+2+k}\int(\dots)r^{-(2+k)}$
is correct for $z_e\neq0$ but $0\cdot\infty$ at $z_e=0$, with a direction-dependent limit: the average convention
of the previous subsection is \emph{wrong} for it when $k\ge1$ ($\int_{\rm half\ circle}\sin\theta\neq\frac12\oint\sin\theta$).
The recursion \eqref{eq:recursion} has no such problem. \status{N}
\end{remark}

\subsection{Half plane and the planar closed form}
\begin{proposition}[Consistency with \cref{sec:planar2}]\label{prop:halfplane}
For $T$ a half plane (a single infinite edge, $z=-d$, indicator $0$, chord $[-\ell,\ell]$, $\ell=\sqrt{1-d^2}$),
\cref{eq:value} equals $\lambda_2(d)$ for every kernel.
\end{proposition}
\begin{proof}
Both sides vanish at $d=1$, so it suffices to compare derivatives in $d$. For the edge side, \cref{thm:grad}
(moving $\xx$ away from the wall by $\dd d$ is a displacement along $\nn$) gives
$\partial_d\lambda=-\int_{-\ell}^{\ell}W(\sqrt{s^2+d^2})\,\dd s=-2\int_d^1W(q)\,q\,(q^2-d^2)^{-1/2}\dd q$ (substituting $q=\sqrt{s^2+d^2}$, $\dd s=q\,\dd q/s$),
which is $\lambda_2'(d)$ by \cref{prop:byparts}.
\end{proof}
The two derivations are independent: the shell calculus integrates arcs, the edge calculus uses the divergence theorem. Note what is and is not
tested: the proof compares the \emph{gradient} identity with the arc form; the value identity itself enters through the single infinite edge
(whose total angle $-\pi$ is closed by the arc at infinity, flux $+\pi M_W(1)$) and is confirmed by the check at $d\to0$ above and
numerically for all four kernels at $d\in\{0,10^{-12},\dots,1-10^{-12}\}$ and symbolically \status{M, N}.

\section{Fields on elements: FEM nodal weights}\label{sec:fem}

Fluid--solid interactions weight a field $A$ living on the solid. On a triangle $T$ let
$A(\xx')=\sum_iA_iN_i(\xx')$ with Lagrange basis $N_i$ of degree $p$.

\begin{proposition}[Nodal weights]\label{prop:fem}
Expand $N_i(\xx+\yy)=\sum_{|\alpha|\le p}B_{\alpha i}(\xx)\,\yy^\alpha$, $B_{\alpha i}=\partial^\alpha N_i(\xx)/\alpha!$ (exact, finite). Then
\begin{align}
 \int_TA\,W\,\dd\xx'&=\sum_iA_iw_i,&& w_i=\sum_\alpha B_{\alpha i}\,m_\alpha,\\
 \int_TA\,\nabla_{\xx} W\,\dd\xx'&=\sum_iA_iG_i,&& G_i=\sum_\alpha B_{\alpha i}\,g_\alpha,
\end{align}
with $g_{\alpha}=\partial_{\xx} m_\alpha+\sum_j\alpha_j\ee_j\,m_{\alpha-\ee_j}$ (the sum runs over the components $j$; $m_{\alpha-\ee_j}=0$ if $\alpha_j=0$), where $\partial_{\xx} m_\alpha$ is
given by \cref{cor:fullspace}.
\end{proposition}
\begin{proof}
The first two are linearity. For $g_\alpha$: $m_\alpha=\int_T\yy^\alpha W(|\yy|)$ with $\yy=\xx'-\xx$, so
$\partial_{x_j}m_\alpha=-\alpha_jm_{\alpha-\ee_j}+g_{\alpha,j}$, because $\partial_{x_j}$ acts on both $\yy^\alpha$ and on $W$,
whereas $g_\alpha$ differentiates $W$ only.
\end{proof}
Hence the weights need exactly the moments $m_\alpha$, $|\alpha|\le p$ ($p\le3$ gives $k\le3$ for $w_i$, and $k\le3$ edge
integrals for $G_i$), computed by \cref{thm:recursion}. \status{N}

\paragraph{Special cases.} $p=0$: $w=m_0$. $p=1$: $N_i$ is the barycentric coordinate, $N_i(\xx+\yy)=N_i(\xx)+\nabla N_i\cdot\yy$, so $w_i=N_i(\xx)\,m_0+\nabla N_i\cdot m_1$.
Partition of unity $\sum_iw_i=\int_TW$ follows from $\sum_iN_i\equiv1$. If $T$ covers the support, the
weights reproduce polynomial fields of degree $\le p$ exactly ($\int A\,W$ and $\int A\nabla_{\xx} W=\int\nabla A\,W$). \status{N}

\paragraph{Conditioning.} The formulas are exact; their floating-point behaviour is not. Three exact
reformulations remove the loss: (1)~for polygons inside the support one may use the potential
$\tilde\Phi=\int_0^rtf$ in \eqref{eq:recursion}, because it differs from $\Phi$ by the constant $C=M_f(r_c)$ (\cref{thm:value}), whose total contribution
$C\big(\oint n_i\yy^\beta\dd s-\beta_i\int_T\yy^{\beta-\ee_i}\big)$ vanishes by the divergence theorem; (2)~the same for $g_\alpha$, which can be written $g_{\alpha,j}=-\int_T\yy^{\alpha+\ee_j}L$ with $L=W'/r$
(because $\partial_{x_j}W=-L\,y_j$); the potential of $L$ is $\Phi[L]=W$, and the shifted potential $\tilde\Phi=W-W(0)$ plays the role of $\tilde\Phi$ in (1); (3)~for far elements the integrand $N_iW$ is smooth and Gauss quadrature converges
spectrally. The practical error tables are in the project notes and tests and are not reproduced here. \status{N}

\section{Closed boundaries (data on the boundary only)}\label{sec:closed}

A solid bounded by a closed polyline is a simple polygon $S$. \Cref{thm:value,thm:grad,thm:recursion} hold
for $S$ unchanged, using only $\partial S$ and the indicator $\ind[\xx\in S]$ (the crossing number, computed from the same
$z_e$ as the edge terms). For a triangulation of $S$ the interior edges cancel pairwise (each appears twice with opposite
$\nn_e$). A field known only on $\partial S$ enters through a polynomial extension of degree $q$ (moments $k\le q$).

\subsection{Polygonal approximation of a curved boundary}
Let $S$ be an inscribed (or circumscribed) regular $N$-gon of a disk of radius $R$, $a=\pi/N$.
In a polar angle $\varphi\in[-a,a]$ about an edge normal (centre of the disk at the origin), the edge sits at radius $\rho_S(\varphi)=R\cos a/\cos\varphi$ (inscribed)
or $R/\cos\varphi$ (circumscribed), so the radial deviation from the circle is
\[
 \delta_{\rm in}(\varphi)=R\Big(\frac{\cos a}{\cos\varphi}-1\Big)=\frac R2(\varphi^2-a^2)+O(Ra^4),\qquad
 \delta_{\rm circ}(\varphi)=R\Big(\frac1{\cos\varphi}-1\Big)=\frac R2\varphi^2+O(Ra^4),
\]
with means over $[-a,a]$ of $-Ra^2/3$ and $+Ra^2/6$ (ratio $-2:1$). The same ratio holds for the area deficit and excess:
$\frac12NR^2\sin2a=\pi R^2(1-\frac23a^2+\dots)$ and $NR^2\tan a=\pi R^2(1+\frac13a^2+\dots)$.
\begin{proposition}\label{prop:ngon}
Let $\epsilon=\frac{Ra^2}{6}\int_0^{2\pi}w\,\dd\phi$ with $w(\phi)=R\,W(|\xx'(R,\phi)-\xx|)$ (the kernel on the circle). If $W\in C^2$ then
$\lambda_{\rm in}-\lambda=-2\epsilon+O(N^{-4})$ and $\lambda_{\rm circ}-\lambda=+\epsilon+O(N^{-4})$, so
$\tfrac13\lambda_{\rm in}+\tfrac23\lambda_{\rm circ}$ is accurate to $O(N^{-4})$. \status{A, N}
\end{proposition}
\begin{proof}[Argument]
In polar coordinates $(\rho,\phi)$ about the disk centre, $S$ and the disk differ in the thin signed ring between $\rho=R$ and $\rho=R+\delta(\phi)$:
\[
 \lambda_S-\lambda=\int_0^{2\pi}\!\!\int_R^{R+\delta}W\rho\,\dd\rho\,\dd\phi=\int_0^{2\pi}w(\phi)\,\delta(\phi)\,\dd\phi+O(R^2a^4).
\]
The deviation $\delta$ is $2a$-periodic with mean $\bar\delta$. Split $\delta=\bar\delta+(\delta-\bar\delta)$. The oscillating part has zero mean on every period, so
two integrations by parts against the $2\pi$-periodic weight $w$ (with $w''$ integrable, which needs $W\in C^2$) give $O(a\cdot a\cdot Ra^2)=O(a^4)$.
Hence $\lambda_S-\lambda=\bar\delta\int w\,\dd\phi+O(a^4)$ with $\bar\delta_{\rm in}=-Ra^2/3+O(a^4)$, $\bar\delta_{\rm circ}=Ra^2/6+O(a^4)$.
\end{proof}
Numerically (w4, $R=\frac32$, $d=0.3$, particle on an edge normal; \path{scripts/derivation_checks/paper_ngon_orders.py}) the errors of both polygons fall by $4$ per
doubling of $N$ and their ratio tends to $-2$. The combination divided by $a^4$ is $0.352,\,0.374,\,0.377,\,0.377,\,0.377$ for $N=24,32,48,64,96$, i.e.\ a clean $O(N^{-4})$ with
coefficient $\approx0.377$ ($1$M and $4$M rays agree to all digits shown). At small $N$ the next term matters: the combination changes sign between $N=12$ ($-0.35\,a^4$) and
$N=24$, passing through zero near $N=16$, which is why a measurement of the order from $N=8,16,32$ alone looked erratic (the earlier project value was $\ge3.3$). \status{N}

\subsection{Limits}
Self-intersecting or overlapping solids double-count (the indicator is a crossing number per solid) and are not covered.
The polygon error $O(N^{-2})$ is exactly the regime where the curvature expansion of \cref{sec:curv} is cheaper than adding edges.

\section{Curvature expansion of the value integral}\label{sec:curv}

For smooth boundaries with radius of curvature large compared to $h$ we expand $\lambda$ in $\kappa$, the
curvature at the closest boundary point, instead of adding edges. Set $h=1$ and take $\kappa>0$ for a locally convex solid
(a disk of radius $R$ has $\kappa=1/R$) and $\kappa<0$ for a concave one (the solid is the complement of a disk).
$\varphi(\sigma)=W(\sqrt\sigma)$, $\sigma=|\yy|^2$, extended by $0$ for $\sigma>1$. The idea: in the coordinates
(arclength $t$ of the closest boundary point, depth $s$) the curved solid is a flat strip with the weight $1-\kappa s$,
and the squared distance to the particle differs from the flat one by a power series in $\kappa$; Taylor-expand $\varphi$ in it.

\subsection{Statement}
\begin{theorem}[Second-order closest-point expansion]\label{thm:curv}
For a disk (more generally, a boundary whose curvature is constant over the support) and a particle at distance $d>0$ outside,
\begin{equation}\label{eq:curv}
 \lambda(d,\kappa)=F_0(d)+\kappa F_1(d)+\kappa^2F_2(d)+O(\kappa^3),
\end{equation}
where, with $y_1=t$, $y_2=d+s$ and integrals over the half plane $\{y_2\ge d\}$,
\begin{align}
 F_0&=\int\varphi,\qquad F_1=\int\Big[\varphi'\,y_1^2\,(2d-y_2)-(y_2-d)\,\varphi\Big],\\
 F_2&=\int\Big[\varphi'\big(-d(y_2-d)y_1^2-\tfrac{y_1^4}{12}\big)+\tfrac12\varphi''y_1^4(2d-y_2)^2-(y_2-d)\varphi'y_1^2(2d-y_2)\Big].
\end{align}
$F_0=\lambda_2$ of \cref{sec:planar2}. The gradient is $\nabla_{\xx}\lambda=\partial_d\lambda\,\hat\nn$ (outward normal) in the constant-curvature case.
\end{theorem}
\begin{proof}
\emph{Step 1 (coordinates).} Write a solid point as $\xx'=c(t)-s\,\nn(t)$, $s$ the depth into the solid and $t$ arclength of the closest
boundary point. For the disk the polar area element $\rho'\dd\rho'\dd\theta$ with $\rho'=R-s$, $\theta=\kappa t$ gives
$\dd\xx'=(1-\kappa s)\,\dd t\,\dd s$. The same holds for $\kappa<0$ ($\rho'=|R|+s$).\\
\emph{Step 2 (distance).} The particle is at distance $R+d$ from the centre, so by the law of cosines
\begin{align*}
 |\yy|^2&=(d+s)^2+(1+\kappa d)(1-\kappa s)\,\frac{2(1-\cos\kappa t)}{\kappa^2}\\
 &=\underbrace{(d+s)^2+t^2}_{\sigma_0}+\kappa\,t^2(d-s)+\kappa^2\Big(-d\,s\,t^2-\frac{t^4}{12}\Big)+O(\kappa^3),
\end{align*}
using $2(1-\cos\kappa t)/\kappa^2=t^2-\kappa^2t^4/12+O(\kappa^4)$ and $(1+\kappa d)(1-\kappa s)=1+\kappa(d-s)-\kappa^2ds$.
\emph{Step 3 (Taylor).} With $\delta_1=t^2(d-s)$ and $\delta_2=-dst^2-t^4/12$, expand
$(1-\kappa s)\,\varphi(\sigma_0+\kappa\delta_1+\kappa^2\delta_2)$ to $O(\kappa^2)$:
\[
 \kappa^1:\ \varphi'\delta_1-s\,\varphi,\qquad
 \kappa^2:\ \varphi'\delta_2+\tfrac12\varphi''\delta_1^2-s\,\varphi'\delta_1 .
\]
These are exactly the integrands of $F_1$ and $F_2$ (with $d-s=2d-y_2$ and $s=y_2-d$).\\
\emph{Step 4 (remainder).} Step~3 is the Taylor expansion of $\varphi$ about $\sigma_0$ with increment
$\delta\sigma=\kappa\delta_1+\kappa^2\delta_2+O(\kappa^3)=O(\kappa)$ on the (bounded) support. For $d\ge d_0>0$ one has $\sigma_0\ge d_0^2$, so $\sigma\ge d_0^2-O(\kappa)$ stays away from
$\sigma=0$, where $\varphi''$ is singular for some kernels ($\varphi''\sim15/r$ for w2 and $\frac94/r$ for the cubic; it is bounded for w4 and w6). Away from $0$, $\varphi''$ is
continuous and piecewise $C^1$, hence Lipschitz, for all four kernels: the cubic has $W'''$ jumping at $q=1$ and $q=\frac12$, but $\varphi''$ involves
only $W$, $W'$ and $W''$, which are continuous. The Taylor remainder after the quadratic term is therefore bounded by $\frac16\mathrm{Lip}(\varphi'')\,|\delta\sigma|^3=O(\kappa^3)$,
uniformly on $d\ge d_0$. (The constant degrades as $d_0\to0$ because of the singularity; this is why the limit $d\to0$ needs separate treatment.)
\end{proof}
Steps 1--3 are verified symbolically \status{M, maple/14 and maple/20}.

\subsection{Evaluation}
Each integral is a sum of half-plane \emph{moments} $m_{ij}[f]=\int_{y_2\ge d}y_1^iy_2^jf(r)$ of the exact piecewise polynomial profiles $W$,
$\varphi'=W'/(2r)$ and $\varphi''=(W'/r)'/(4r)$. As polynomials in $r$ these reach down to $r^{-1}$ for w2 (where $\varphi''=15/r+\dots$), to $r^{0}$ for w4 and w6,
and for the cubic to $r^{-3}$ in each block (the singular parts of the two blocks cancel to a net $\frac94/r$). Each moment is a single edge
instance of \cref{thm:recursion} (here $i+j\le6$; the recursion is general, its symbolic check in the project covers $k\le4$): no new quadrature. In practice $F_k(d)$ are tabulated once as one-dimensional functions.

\subsection{Accuracy}
The measured errors (Wendland C4, $d=0.3h$) follow the predicted orders:
\begin{center}
\begin{tabular}{rrrrr}
\toprule
$\kappa h$ & exact $\lambda$ & planar $F_0$ & $+\kappa F_1$ & $+\kappa^2F_2$\\
\midrule
1/2 & 0.09017 & $9.0\!\cdot\!10^{-3}$ & $9.9\!\cdot\!10^{-4}$ & $1.4\!\cdot\!10^{-4}$\\
1/4 & 0.09444 & $4.7\!\cdot\!10^{-3}$ & $2.6\!\cdot\!10^{-4}$ & $1.8\!\cdot\!10^{-5}$\\
1/8 & 0.09674 & $2.4\!\cdot\!10^{-3}$ & $6.8\!\cdot\!10^{-5}$ & $2.3\!\cdot\!10^{-6}$\\
1/16& 0.09794 & $1.2\!\cdot\!10^{-3}$ & $1.7\!\cdot\!10^{-5}$ & $3.0\!\cdot\!10^{-7}$\\
\bottomrule
\end{tabular}
\end{center}
(errors against the exact disk, \status{N}; recomputed independently with a shell quadrature of the exact disk and the half-plane integrals $F_0,F_1,F_2=0.09917,\,-0.01999,\,0.004497$, all entries reproduced);
successive halving of $\kappa$ divides the three error columns by $\approx2,4,8$.

\subsection{General curves}
\begin{remark}[Variable curvature]\label{rem:varcurv}
Let the boundary have $\kappa(t)=\kappa+\kappa't+\tfrac12\kappa''t^2+\dots$ about the closest point of the particle. Count orders by dimension:
$\kappa\sim h^{-1}$, $\kappa'\sim h^{-2}$, $\kappa''\sim h^{-3}$. Thus $\kappa'$ could enter \cref{eq:curv} at order $2$, next to $\kappa^2$.
It does not: the particle lies on the normal through the closest point, so the reflection $t\mapsto-t$ maps the particle to itself, sends $\kappa'\mapsto-\kappa'$ and
leaves $\lambda$ unchanged. Hence $\lambda$ is even in $\kappa'$ and has no term linear in $\kappa'$; the first correction beyond \cref{eq:curv} is $\kappa''$ (and
$\kappa'^2$ at order $4$), both of order $\ge3$. \status{N}
The \emph{gradient} is different. With $\kappa(\xx)$ the curvature at the closest boundary point, $\nabla\lambda=\partial_d\lambda\,\hat\nn+F_1\nabla\kappa+\dots$
($\kappa$ is constant along normals), so there is a tangential part $\kappa'F_1\,\hat{\bm t}$ of order $2$. The gradient formula of \cref{thm:curv}
therefore has an $O(\kappa^2)$ tangential error unless $\kappa'F_1$ is added; for constant curvature it is exact to the order of the value expansion. At the next order the
tangential component is $\kappa'\big[F_1+\kappa(2F_2-dF_1)\big]$: the factor $2F_2$ is $\partial_\kappa(\kappa^2F_2)$, and $-dF_1$ comes from $\partial_x\kappa=\kappa'/(1+\kappa d)$ (a tangential
displacement $\varepsilon$ of the particle moves the closest point by $\varepsilon/(1+\kappa d)$). \status{N}
Numerically (w4, $d=0.3$, boundary $y=-\frac\kappa2x^2-\frac{\kappa'}6x^3$, \path{scripts/derivation_checks/paper_curvature_variable.py}): $\lambda(\kappa,\kappa')-\lambda(\kappa,0)$ falls by $4$ when $\kappa'$ is halved (even in $\kappa'$);
the parabola with $\kappa'=\kappa^2$ is reproduced by $F_0+\kappa F_1+\kappa^2F_2$ with error falling by $8.0$ per halving of $\kappa$; the tangential gradient at $(\kappa,\kappa')=(\frac18,\frac1{64})$ is
$-2.840\cdot10^{-4}$ against $-3.123\cdot10^{-4}$ for $\kappa'F_1$ alone and $-2.830\cdot10^{-4}$ with the correction.
\end{remark}
\begin{remark}[Validity]
The expansion is asymptotic in $\kappa h$ (usable for $\kappa h\lesssim\frac14$ at $10^{-5}$), requires reach $\ge h$ on the solid side, a single
boundary sheet in the support, and $R>h-d$ so the support does not engulf the disk; the limit $d\to0$ needs the limit form of the negative-power primitives.
\end{remark}

\section{Slender and thin obstacles}\label{sec:slender}

When a body is thin compared to $h$ (fibres, strands) its volume integral reduces to a line integral with a short correction series.

\begin{lemma}[Pizzetti mean-value series]\label{lem:pizzetti}
For $f$ real-analytic on $\overline{B_a(c)}\subset\R^2$,
\[
 \frac1{\pi a^2}\int_{B_a(c)}f=\sum_{k\ge0}\frac{(a^2/4)^k}{k!\,(k+1)!}\,\Lap^kf(c).
\]
\end{lemma}
\begin{proof}
Let $u(\rho)=\frac1{2\pi}\oint f(c+\rho\bm\omega)\,\dd\bm\omega$ be the circle mean. Averaging commutes with $\Lap$, and for a radial
function the 2D Laplacian is $u''+u'/\rho$, so $u''+u'/\rho=\overline{\Lap f}(\rho)$, the circle mean of $\Lap f$ (Darboux).
Insert $u=\sum_ku_k\rho^{2k}$ with $u_0=f(c)$: $(2k)^2u_k=\big[\rho^{2k-2}\big]\overline{\Lap f}=u_{k-1}[\Lap f]$, so by induction $u_k=\Lap^kf(c)/(4^k(k!)^2)$
(this is Pizzetti's formula~\cite{pizzetti1909}). Average over $\rho\in[0,a]$ with weight $2\rho/a^2$: $\int_0^a2\rho^{2k+1}\dd\rho/a^2=a^{2k}/(k+1)$, which turns $\frac{1}{4^k(k!)^2}$ into $\frac{1}{4^kk!(k+1)!}$.
\end{proof}

\begin{proposition}[Disk in 2D, strand in 3D]\label{prop:slender}
(a) A disk of radius $a$ at distance $D$ from the particle, with $W$ smooth on $[D-a,D+a]$:
$\int_{B_a}W=\pi a^2\sum_{k\le K}\frac{(a^2/4)^k}{k!(k+1)!}\Lap^kW(D)+O(a^{2K+4})$.\\
(b) Straight 3D strand of radius $a$ along a segment: with $\Lap_\perp$ the Laplacian in the cross-section,
\[
 \int_{\rm strand}W\simeq\pi a^2\int_{\rm axis}\Big[W+\frac{a^2}{8}\Lap_\perp W+\frac{a^4}{192}\Lap_\perp^2W+\dots\Big]\dd s,
\]
the coefficients $\frac1{4\cdot1!\,2!}$ and $\frac1{16\cdot2!\,3!}$ of \cref{lem:pizzetti}.\\
(c) 2D strip of half width $a$: $\int_{-a}^a f(z+u)\dd u=\sum_k\frac{2a^{2k+1}}{(2k+1)!}f^{(2k)}(z)$.
\end{proposition}
\begin{proof}
(a),(b) apply \cref{lem:pizzetti} (in the cross-section for (b)); the first neglected term is $\pi a^2\cdot a^{2K+2}$. (c) is the Taylor series integrated
over $u$. The leading term of (b) is the line integral of $W$ along the axis chord, $\pi a^2\int_{\rm chord}W\,\dd s$ with $z$ the distance of the particle to the axis. \status{M for the coefficients and for (c);
N for the orders $4,6,8$ of (a), the 2D case: error ratios $15.7$, $63.6$, $256.1$ per halving of $a$ for $K=0,1,2$ (w4, $D=0.5$, \path{paper_curvature_variable.py})}
\end{proof}

\begin{remark}[Regularity]
Wendland kernels have no linear term, so $\Lap_\perp W$ is regular on the axis; but $C^2$ Wendland has an $r^3$ term, hence $\Lap_\perp^2W\sim1/r$
is singular where the axis passes through the particle, and the series is asymptotic and breaks at kernel knots or the support rim. Where it fails the
polygon of \cref{sec:closed} is the fallback; no gap exists in 2D since the strip is exactly $\lambda_2(d-a)-\lambda_2(d+a)$.
\end{remark}

\paragraph{Ball in 3D.} For a ball of radius $\rho$ at centre distance $D$ the \emph{shell theorem} reads
$\int_{S^2}f(|\xx-\bm c-\rho\bm\omega|)\,\dd\bm\omega=\frac{2\pi}{D\rho}\int_{|D-\rho|}^{D+\rho}r\,f(r)\,\dd r$ \status{M}, consistent with \cref{prop:sphere}.

\paragraph{Curved centre line.} A first-order correction $\propto\kappa_ca^2$ is expected (the tube volume is independent of $\kappa_c$ but the kernel weighting is not
symmetric); not derived. \status{P}

\section{Extension to three dimensions}\label{sec:threed}

\begin{theorem}[Face identity for a polyhedron $T$]\label{thm:face}
For $g$ radial with support $R$ and $M_3(r)=\int_0^rt^2g\,\dd t$,
\[
 \int_Tg=\ind[\xx\in T]\,4\pi M_3(R)+\sum_{f\in\partial T}z_f\int_f\frac{M_3(r)-M_3(R)}{r^3}\,\dd A,
\]
$z_f=\nn_f\cdot(p_f-\xx)$. For a normalised kernel $4\pi M_3(1)=1$.
\end{theorem}
\begin{proof}
$\nabla\!\cdot(\yy M_3/r^3)=3M_3/r^3+r(M_3'/r^3-3M_3/r^4)=M_3'/r^2=g$ [M]. Apply the divergence theorem on $T\setminus B_\varepsilon(\xx)$
(small-sphere flux $-4\pi M_3(\varepsilon)\to0$). For the constant part, $\sum_fz_f\int_f\dd A/r^3$ is the total solid angle subtended by $\partial T$
($=4\pi$ inside, $0$ outside).
\end{proof}
The integrand on each face is compact: it is a 2D value-type integral in the face plane with profile $z_f(M_3(r)-M_3(R))/r^3$,
$r^2=\rho^2+z_f^2$, to which \cref{thm:value} (in the plane, with the foot of $\xx$) applies, leaving edge integrals of the form
$\int r^m/(s^2+w_e^2)\,\dd s$ with $r^2=s^2+w_e^2+z_f^2$, which are elementary (solid-angle-type $\arctan(sz/(wr))$,
cf.\ Van Oosterom--Strackee~\cite{vanoosterom1983}).

\paragraph{Status.} The identity above is proved; the \emph{evaluation} (edge primitives in 3D, moment recursion in the face plane, comparison with the
planar $\lambda_3$ and the sphere closed form, polyhedral potentials) is a plan. \status{P}
The divergence step is verified symbolically; nothing else in this section has been numerically checked.

\section{Wall operators built on the edge calculus}\label{sec:wall}

The same edge integrals provide terms of the SPH momentum equation. Here $\yy=\xx'-\xx$ and $\xx$ is a fluid particle,
the solid half plane being locally the wall. Quantities are 2D. The solid is modelled as a continuum of mass $m_{\rm wall}$ per unit area, measured in units of
the reference density $\rho_0$, so a sum $\sum_jV_j(\cdot)$ over solid particles of volume $V_j=m_j/\rho_j$ becomes $\frac{m_{\rm wall}}{\rho_i}\int_S(\cdot)\,\dd\xx'$ with $\rho_i$ the density
of the fluid particle (also in units of $\rho_0$).

\subsection{The wall Laplacian}
\begin{proposition}[Green form]\label{prop:lapwall}
For the solid polygon $S$ define $\lambda_\Delta(\xx):=\int_S\nabla^2_{\xx}W(|\xx-\xx'|)\,\dd\xx'$. Then
\begin{equation}\label{eq:lapwall}
 \lambda_\Delta=\sum_ez_e\int_{\chord_e}\frac{W'(r)}{r}\,\dd s.
\end{equation}
For a fluid-side particle $z_e<0$ and $W'<0$, so $\lambda_\Delta>0$; $\lambda_\Delta=0$ when the support lies inside $S$ or is disjoint from it.
\end{proposition}
\begin{proof}
$\nabla_{\xx}^2W(|\xx-\xx'|)=\nabla_{\xx'}^2W$, so Green's theorem gives $\oint_{\partial S}\partial_nW\,\dd s$ with $\nabla_{\xx'}W=W'(r)\yy/r$ and $\nn_e\!\cdot\!\yy=z_e$.
If $\xx\in S$ excise a small disk: its boundary flux is $2\pi\varepsilon\,W'(\varepsilon)\to0$ since $W'(r)=O(r)$ at the centre, so there is no angle term;
the outer support circle contributes nothing because $W'(1)=0$. If the support lies inside $S$ the boundary is outside the support, so the sum is empty; if disjoint, the chords are empty.
\end{proof}
Numerically, \cref{eq:lapwall} agrees with a ray quadrature of $\int_S\nabla^2W$ over a non-convex polygon to $\le8\cdot10^{-11}$ for all four kernels, for a particle inside, outside and near an edge
(\path{paper_edge_identities_k6.py}). \status{N}

\paragraph{Evaluation with the existing operators.} Let $L=W'/r$. Then $W'=rL$, $W''=L+rL'$ and $\nabla^2W=W''+W'/r=2L+rL'$ in 2D [M]. With the density $\lambda[L]=\int_SL$ and
the covariance $\mathrm{Cov}[L]=\int_S\yy\otimes\nabla_{\xx} L$ (where $\nabla_{\xx} L=-L'\yy/r$ for $\yy=\xx'-\xx$, hence $\operatorname{tr}\mathrm{Cov}[L]=-\int_SrL'$),
\begin{equation}\label{eq:lapscene}
 \lambda_\Delta=2\lambda[L]-\operatorname{tr}\mathrm{Cov}[L].
\end{equation}
Registering $L$ rather than $\nabla^2W$ as the kernel is essential: inside a body the indicator contribution is $2\cdot1-\operatorname{tr}(\mathbb 1)=0$ and cancels
(over all of $\R^2$, $\mathrm{Cov}=\mathbb 1\int L$ by parts, so $2\int L-2\int L=0$, as $\int\nabla^2W=0$ demands),
whereas a directly registered $\nabla^2W$ would give a spurious constant: the scene assumes $\int K=1$ for a registered kernel, but $\int\nabla^2W=0$.

\paragraph{Wendland kernels: $L$ is a multiple of an ordinary kernel.} With $W=C_2h^{-2}\what(q)$ one finds $W'/r=C_2h^{-4}\,\what'(q)/q$ and
\[
 \what_{\rm w2}'/q=-20(1-q)^3,\qquad \what_{\rm w4}'/q=-\tfrac{56}3\,(1+5q)(1-q)^5=-\tfrac{56}3\,\big(6(1-q)^5-5(1-q)^6\big)\quad\text{[M]}.
\]
Normalise the shapes $(1-q)^3$ and $6(1-q)^5-5(1-q)^6$ as 2D kernels $K_\ell=C_\ell\,\mathrm{shape}(q)/h^2$ with $2\pi C_\ell\int_0^1q\,\mathrm{shape}\,\dd q=1$:
since $\int_0^1q(1-q)^3\dd q=\frac1{20}$ and $\int_0^1q(1+5q)(1-q)^5\dd q=\frac3{56}$, $C_\ell=\frac{10}\pi$ (w2) and $\frac{28}{3\pi}$ (w4). Then $L=f\,K_\ell$ with
\[
 f=\frac{P\,C_2}{C_\ell\,h^2}=-\frac{14}{h^2}\ \text{(w2: }P=-20,\ C_2=\tfrac7\pi),\qquad f=-\frac{18}{h^2}\ \text{(w4: }P=-\tfrac{56}3,\ C_2=\tfrac9\pi),
\]
where $P$ is the constant in front of the shape in the display above. \status{M}

\subsection{The viscosity coefficient}
\begin{lemma}\label{lem:nu}
For every normalised radial kernel in $n$ dimensions, $\int_{\R^n}r\,W'(r)\,\dd\xx=-n$; for $n=2$, $-2$.
\end{lemma}
\begin{proof}
In 2D, $2\pi\int_0^1r^2W'\dd r=2\pi\big([r^2W]_0^1-2\int_0^1rW\dd r\big)=-2$, since $2\pi\int rW\,\dd r=1$ is the normalisation. (Likewise in $n$ dimensions by one integration by parts.) \status{M}
\end{proof}

The solver's artificial viscosity is, for fluid pairs, the pairwise sum (with $\mathrm{fac}=\alpha c_0h/\xi$ the combination of the scheme's constants)
\begin{equation}\label{eq:pairwise}
 \bm a_i=\mathrm{fac}\sum_j\frac{m_j}{\bar\rho_{ij}}\,\frac{(\bm v_i-\bm v_j)\cdot(\xx_i-\xx_j)}{|\xx_i-\xx_j|^2}\,\nabla_iW_{ij},
\end{equation}
which we take from the code (\texttt{deltasph2d.py}) as given. Its continuum limit fixes the coefficient $\nu_{\rm eff}$:

\begin{proposition}[Pairwise bulk term]\label{prop:nueff}
For a smooth velocity field the sum \cref{eq:pairwise} converges to $\bm a=\dfrac{\mathrm{fac}}{2(n+2)}\big(\nabla^2\bm v+2\nabla\nabla\!\cdot\bm v\big)$; in 2D $\nu_{\rm eff}=\mathrm{fac}/8=\alpha c_0h/(8\xi)$ is the coefficient of $\nabla^2\bm v$. \status{M (moments), N (lattice), A (form \cref{eq:pairwise})}
\end{proposition}
\begin{proof}
Replace $\sum_jm_j/\bar\rho_{ij}$ by $\int\dd\xx'$, put $\yy=\xx'-\xx$, $\delta\bm v=\bm v(\xx')-\bm v(\xx)$. Then $\bm v_i-\bm v_j=-\delta\bm v$, $\xx_i-\xx_j=-\yy$ and $\nabla_iW_{ij}=-W'\yy/r$, so
$a_i=-\mathrm{fac}\int\frac{W'}{r^3}\,y_i\,y_k\,\delta v_k\,\dd\xx'$. Taylor-expand $\delta v_k=y_l\partial_lv_k+\frac12y_ly_m\partial_l\partial_mv_k+\dots$.
Odd moments vanish by symmetry. For the quadratic term the isotropic fourth moment
\[
 \int y_iy_ky_ly_m\,F(r)\,\dd\xx=\frac{\delta_{ik}\delta_{lm}+\delta_{il}\delta_{km}+\delta_{im}\delta_{kl}}{n(n+2)}\int r^4F\,\dd\xx
\]
(check by contracting $i=k$, $l=m$) with $F=W'/r^3$ gives $\int r^4F=\int rW'=-n$ by \cref{lem:nu}. Hence
$a_i=-\mathrm{fac}\cdot\frac12\cdot\frac{-n}{n(n+2)}\big(\partial_l\partial_lv_i+2\partial_i\partial_kv_k\big)$, which is the claim.
\end{proof}
Numerically, for generic quadratic velocity fields and all four kernels the continuum integral reproduces $\frac18(\nabla^2\bm v+2\nabla\nabla\!\cdot\bm v)$ to $10^{-13}$ (cubic $2\cdot10^{-9}$, quadrature of the knot) and the lattice sum \cref{eq:pairwise} with $h/\Delta x=16$ to $\le1.2\cdot10^{-5}$ (\path{scripts/derivation_checks/paper_viscosity_limits.py}). The only moment of the kernel that survives is $\int rW'$, which is fixed by the normalisation: the coefficient does not depend on the kernel family.
Note that the bulk term contains the extra part $2\nabla\nabla\!\cdot\bm v$; the wall terms below carry only the $\nabla^2$ part, so wall and bulk share the Laplacian coefficient but are not one operator.

\subsection{Free-slip wall term}
A Laplacian estimate of the smoothed field, exact for the convolved field, is $\nabla^2\bm v(\xx)\simeq\int(\bm v(\xx')-\bm v(\xx))\,\nabla^2W(|\xx-\xx'|)\,\dd\xx'$ (because $\int\nabla^2W=0$).
(Equivalently, by $\int rW'=-n$, $\nabla^2\bm v=-2\int(\bm v'-\bm v)\,W'/r$, the Morris form; both estimators converge as $O(h^2)$, error ratios $4.0$ per halving of $h$ for w2, w4 and the cubic.) The solid part of the integral needs ghost values. With the free-slip mirror, constant normal $\nn$ per particle (pointing \emph{into} the wall) and
$u_n=(\bm v-\bm v_w)\cdot\nn$, the ghost velocity is the mirror image, $\bm v_{\rm ghost}-\bm v=-2u_n\nn$, constant over the solid. Then
\begin{equation}\label{eq:freeslip}
 \bm a_{\rm wall}=\nu_{\rm eff}\,\frac{m_{\rm wall}}{\rho_i}\int_S(\bm v_{\rm ghost}-\bm v)\,\nabla^2W\,\dd\xx'=-2\nu_{\rm eff}\,\frac{m_{\rm wall}}{\rho_i}\,u_n\,\lambda_\Delta\,\nn .
\end{equation}
Only the normal relative velocity is damped, and it is damped when moving into the wall ($u_n>0$, $\lambda_\Delta>0$). This differs from the pairwise free-slip form near the wall: $\lambda_\Delta\to0$
linearly as the particle reaches the wall (the support loses half its disk and the odd mirror field vanishes), while the pairwise coefficient does not, by design. \status{A}

\subsection{No-slip flux term}
Following Chiron et al.~\cite{chiron2019}, $\int_{\chord}W\,\dd s$ over the cut face is the per-edge channel of $\nabla\lambda$ (\cref{thm:grad}). Define
$|G|=m_{\rm wall}\int_{\chord}W\,\dd s$ and
\begin{equation}
 \bm a_{\rm wall}=-2\nu_{\rm eff}\,(\bm v_i-\bm v_w)\,|G|\,/\,(\rho_i\,d_b),\qquad d_b=\max(d,d_{\min}).
\end{equation}
This is a \emph{modelling} step: it replaces the normal derivative at the wall by the one-sided difference $(\bm v-\bm v_w)/d$ (neglecting the tangential gradient, Chiron et al.\ Eq.~89--90) and uses the viscous coefficient $\nu_{\rm eff}$
of the scheme (an artificial coefficient, about $3000$ times the kinematic viscosity of water in the tested tank). The factor $2$ is the same $2$ as in the Morris form above ($\nabla^2\bm v=2\int(\bm v-\bm v')W'/r$, a consequence of $\int rW'=-n$), which in Chiron's operator multiplies the fluid sum and the wall surface term alike: their Eq.~(92) is
$\frac2{\gamma_i}\big[\sum_j\omega_j(u_i-u_j)\frac{\xx_i-\xx_j}{|\xx_i-\xx_j|^2}\!\cdot\!\nabla_iW_{ij}+\sum_s s_s\frac{u_i-u_s}{d_n}W_{is}\big]$, and in their derivation (Eq.~(84)--(91)) the $2$ in front of the
surface term comes from $\nabla f(\xx)\!\cdot\!\nn\approx\nabla f(\yy)\!\cdot\!\nn$ in the boundary flux $\int_{\partial\Omega}(\nabla f(\xx)\!\cdot\!\nn+\nabla f(\yy)\!\cdot\!\nn)W\,\dd S$. (With $\nn$ outward from the fluid their
$d_n$ is signed so that the wall term opposes $u_i-u_s$, as in the formula above.) None of
this follows from the preceding calculus, so only the identity $|G|=m_{\rm wall}\int W\,\dd s$ is a theorem here. \status{A}

\section{Verification ledger}\label{sec:verification}

\Cref{tab:ledger} lists each result with the strongest evidence behind it.
\begin{table}[h]
\centering\small
\begin{tabular}{p{4.2cm}p{1.3cm}p{6.4cm}}
\toprule
result & tag & evidence\\
\midrule
Kernel constants, block weights & M & \texttt{20\_paper\_audit.mpl}\\
Planar 3D closed forms & M & same (cubic both branches, degrees $8,11,14$)\\
Planar 2D by-parts, $\mathcal J_m$ recursion, w2 explicit form, cubic join & M & same, 40 digits where numeric\\
Sphere cap area, planar limit & M & same\\
Edge primitives & M & \texttt{10\_edge\_primitives.mpl}\\
Blocks, value/gradient, half plane & M,N & \texttt{11,13}, \texttt{tests/edge/test\_value.py}, \texttt{test\_gradient.py}\\
Moment recursion $k\le4$ (symbolic), $k\le6$ (numeric) & M,N & \texttt{12\_moments\_recursion.mpl}; \path{paper_edge_identities_k6.py}\\
FEM weights $p\le3$ & N & \texttt{tests/edge/test\_fem*.py}\\
Polygon-to-disk orders & A,N & argument in \cref{prop:ngon}; \texttt{test\_tier2.py}; \path{paper_ngon_orders.py} ($N=12$--$96$, $\to0.377\,a^4$)\\
Curvature expansion & M,N & \texttt{14\_tier3\_2d.mpl}, \texttt{20\_paper\_audit.mpl}; table reproduced by \path{paper_curvature_table.py}\\
Disk/strip series (2D) & M,N & \texttt{20\_paper\_audit.mpl}, \texttt{test\_tier4.py}, \path{paper_curvature_variable.py}\\
Variable curvature (even in $\kappa'$; tangential gradient) & N & \path{paper_curvature_variable.py}\\
Wall Laplacian Green form & M,N & identities in \texttt{20\_paper\_audit.mpl}; \texttt{tests/scene/test\_viscosity\_scene.py}\\
$\nu_{\rm eff}$ from the pairwise term & M, N, A & \cref{prop:nueff}: moments in \texttt{20\_paper\_audit.mpl}, lattice in \path{paper_viscosity_limits.py}; the pairwise form \cref{eq:pairwise} is quoted from the code\\
Free-slip wall term & A & \cref{eq:freeslip}: ghost-value argument, model\\
No-slip coefficient & A & modelling choice\\
3D face identity & M (div.) & divergence step only\\
3D evaluation, curved centre line & P & open\\
\bottomrule
\end{tabular}
\caption{Verification ledger.}\label{tab:ledger}
\end{table}

Reproduce (all Maple lines must read PASS):
\begin{itemize}
\item \texttt{make -C paper} builds the paper;
\item \texttt{maple -q maple/20\_paper\_audit.mpl}: symbolic and 40-digit checks of this paper (110 PASS lines);
\item \texttt{bash maple/run\_edge.sh}: the project's edge-calculus checks (61 lines at the time of writing);
\item numerical checks in \texttt{scripts/derivation\_checks/} (compare with the quoted numbers):
\begin{itemize}
\item \texttt{paper\_ngon\_orders.py}, \texttt{paper\_curvature\_table.py} (seconds);
\item \texttt{paper\_curvature\_variable.py}, \texttt{paper\_viscosity\_limits.py} (seconds);
\item \texttt{paper\_edge\_identities\_k6.py} (about 4~minutes).
\end{itemize}
\end{itemize}


\appendix
\section{Audit}\label{app:audit}
\subsection{The existing derivation notes}

The mathematical content of \texttt{docs/derivation.md} and \texttt{docs/derivations/*.md} was re-derived and re-checked
(Maple, exact arithmetic or 40 digits). \textbf{No error was found in any final formula.} The following are
presentation or scope issues, in decreasing importance.
\begin{enumerate}
\item \textbf{Sign convention in \texttt{laplacian-wall.md} \S3.} Step~1 defines $\yy=\xx-\xx'$, Step~2 then writes $\operatorname{tr}\mathrm{Cov}=\int rL'$ but the header and
the conclusion need $\operatorname{tr}\mathrm{Cov}=-\int rL'$ (true for $\yy=\xx'-\xx$). The final identity \cref{eq:lapscene} is correct with $\yy=\xx'-\xx$; the intermediate line is not.
\item \textbf{Tier-3 gradient scope} (\texttt{tier3-curvature-2d.md} \S1): ``$\nabla_{\xx}\lambda=\partial_d\lambda\,\hat\nn$'' holds only for constant curvature; for general curves
a tangential term $\kappa'F_1$ appears at second order (\cref{rem:varcurv}). The value expansion needs no $\kappa'$ term.
\item \textbf{$J_m$ structure} (\texttt{derivation.md} \S3.2): odd $m$ are $\mathrm{poly}(d^2)\sqrt{1-d^2}$ (no constant), by \cref{eq:Jrec}.
\item \textbf{Edge value identity, \texttt{tet-face-edge-chain.md}}: status \textsf{[P]} is accurate; the divergence step is nonetheless a theorem (\cref{thm:face}) and could be tagged as such.
\item \textbf{$\nu_{\rm eff}$ and the no-slip factor $2$} (\texttt{laplacian-wall.md} \S3.5, \texttt{noslip-wall.md} \S3.3) are tagged \textsf{[V]} in the notes because the \emph{implementation}
matches the stated formula; as mathematics, the first depends on the pairwise term (not re-derived) and the second is a modelling choice.
\item \textbf{Block identity wording.} \texttt{truncated-monomials.md} correctly notes that block-by-block gradients need no cancellation between blocks; this is \cref{thm:grad}.
The known bug in \texttt{scripts/derivation\_checks/monomial\_edge\_forms.py} (cubic correction $-4\sigma$ instead of $-8\sigma$) remains in the unused script.
\end{enumerate}
Not independently audited here: the solver-side statements in \path{docs/deltasph-*.md} and \path{docs/dfsph-validation.md}, which are validation reports rather than derivations.

\subsection{The first draft of this paper}
Each derivation of the first draft was re-done independently by hand, the symbolic script re-run, and the numerical claims recomputed (the curvature table and
the cubic value at $d=\frac12$ with quadrature, the polygon orders with a ray integration against the disk shell formula). Again no final formula was wrong, but the draft contained the
following errors or gaps, now corrected in the text, in decreasing importance:
\begin{enumerate}
\item \textbf{Remainder in \cref{thm:curv}.} The draft said that $\varphi$ is $C^2$ with $\varphi''$ Lipschitz for all four kernels and argued with a jump set of measure $O(\kappa)$.
$\varphi''$ is in fact singular at $\sigma=0$ ($15/r$ for w2, $\frac94/r$ for the cubic). The result is unaffected because $\sigma\ge d_0^2-O(\kappa)>0$, and where $\varphi''$ is Lipschitz the remainder
bound is immediate, so the measure argument was unnecessary. The profile list ``powers of $r$ down to $r^{-1}$ for Wendland'' was likewise true only for w2.
\item \textbf{Block weights.} The draft called $\frac87$ and $-\frac17$ both the indicator weight and the arctan coefficient; they are the indicator weights $2\pi M_j(r_j)$, and the arctan coefficient is $M_j(r_j)=w_j/2\pi$.
\item \textbf{Edge crossing (\cref{sec:degenerate}).} The draft wrote that the arctan term jumps by $+1$ ``in the sense that cancels'' the indicator jump; in \cref{eq:value} the chord term jumps by $-1$ (the arctan appears with a minus sign in \cref{eq:block}),
which is what cancels. The vertex convention is now justified by an explicit limit.
\item \textbf{Primitives for odd $j$.} ``The same with $J$'' hid that the binomial runs over $p=\lfloor j/2\rfloor$ and the exponent is $m+2i+2$; written out in \cref{prop:primitives}(c).
\item \textbf{Cubic value at $d=\frac12$.} Closed form correct, decimal wrong: $0.03744$, not $0.03748$.
\item \textbf{\Cref{prop:halfplane}.} The draft called the half-plane check a proof of the edge calculus against the shell calculus; the proof tests the gradient identity, and the value identity enters only through the closing at infinity. Reworded.
\item \textbf{$\nu_{\rm eff}$.} Marked ``not re-derived''; it is now derived from the solver's pairwise form (\cref{prop:nueff}). The form itself is still taken from the code. The factor $2$ of the no-slip term is traced to the Morris normalisation instead of ``by analogy'', with the caveat that Chiron's Eq.~92 is quoted from the project notes.
\item \textbf{Scaling lemma.} $m_\alpha(h)=h^{|\alpha|}m_\alpha(1)$ ``after scaling all lengths'' was ambiguous; now $m_\alpha(d;h)=h^{|\alpha|}m_\alpha(d/h;1)$.
\item \textbf{Variable-curvature gradient.} ``First-order accurate'' replaced by the actual statement: an $O(\kappa^2)$ tangential error unless $\kappa'F_1$ is added.
\item \textbf{Polygon orders (\cref{prop:ngon}).} The ``hence proportional to the mean deviation'' step is now an argument (ring integral, two integrations by parts) with the exact hypothesis ($W\in C^2$).
\item \textbf{Variable-curvature gradient (\cref{rem:varcurv}).} The leading tangential term $\kappa'F_1$ is right, but the draft stopped there; measured numerically, the next term is $\kappa\kappa'(2F_2-dF_1)$, with the $-dF_1$ from $\partial_x\kappa=\kappa'/(1+\kappa d)$.
\end{enumerate}
Presentation only: the symbols $R$ (cut radius versus body radius), $\lambda_i$ (barycentric versus kernel integral), $J_m$ (planar family versus edge primitive), $\Phi$ (compact potential versus
kernel profile), $\Delta\lambda$ (read as a difference) and $H$ versus $h$ each had two meanings; see \cref{sec:notation}. $\lambda_\Delta$, $P$, $c$ and $C_\ell$ were also undefined where used.

\subsection{Independent numerical verification of the 2D claims}
The 2D claims that were tagged \textsf{A} or checked only inside the project were re-verified with code that does not use the project library
(\path{scripts/derivation_checks/paper_*.py}): moment recursion for all $|\alpha|\le6$ and gradient, $g_\alpha$ and Green-form identities against ray quadrature
(\cref{sec:edge,sec:wall}); the pairwise-viscosity limit and the two Laplacian estimators (\cref{sec:wall}); the variable-curvature statements and the disk-series orders
(\cref{sec:curv,sec:slender}); the polygon orders over $N=12,\dots,96$ (\cref{sec:closed}). The Chiron et al.\ Eq.~(92) was read in the original. Not covered: everything in
\cref{sec:threed} and the 3D strand, curved centre line and sphere statements; the pairwise form \cref{eq:pairwise} itself, which is quoted from the solver; and the free-slip and no-slip wall terms as models.


\begin{thebibliography}{9}
\bibitem{wendland1995} H.~Wendland, Piecewise polynomial, positive definite and compactly supported
radial functions of minimal degree, \emph{Adv.\ Comput.\ Math.}\ 4 (1995) 389--396.
\bibitem{winchenbach2020} R.~Winchenbach, H.~Akhunov, A.~Kolb, Semi-analytic boundary handling below
particle resolution for smoothed particle hydrodynamics, \emph{ACM Trans.\ Graph.}\ 39(6) (2020).
\bibitem{winchenbach2025} R.~Winchenbach, A.~Kolb, arXiv:2507.21686 (2025), Sec.~6.2 (FEM-weighted
boundary integrals, validation problems).
\bibitem{chiron2019} L.~Chiron et al., SPH wall boundary conditions for three-dimensional complex
geometries (2019), the wall Laplacian of Sec.~6.1 (renormalised Morris operator).
\bibitem{vanoosterom1983} A.~Van Oosterom, J.~Strackee, The solid angle of a plane triangle,
\emph{IEEE Trans.\ Biomed.\ Eng.}\ 30 (1983) 125--126.
\bibitem{pizzetti1909} P.~Pizzetti, Sulla media dei valori che una funzione dei punti dello spazio
assume alla superficie di una sfera, \emph{Rend.\ Lincei} 18 (1909) 182--185.
\end{thebibliography}

\end{document}
